% Bab I  Pendahuluan
\chapter{PENDAHULUAN}

\section{Pengertian}
Skripsi adalah karya tulis ilmiah yang disusun berdasarkan hasil penelitian lapang atau laboratorium yang dilengkapi dengan analisis statistik inferensial untuk pendekatan kuantitatif dan deskriptif untuk pendekatan kualitatif. Pada Program Studi Rekayasa Sistem Komputer, skripsi merupakan tugas akhir mahasiswa program sarjana (S1) sebagai salah satu syarat memperoleh gelar Sarjana Komputer (S.Kom.).

Skripsi disusun secara mandiri di bawah bimbingan dua orang dosen pembimbing. Isinya harus menunjukkan kemampuan mahasiswa dalam mengidentifikasi masalah, mengkaji teori dan penelitian terdahulu, merancang dan melaksanakan penelitian dengan metode ilmiah, menganalisis data, serta menyimpulkan dan melaporkan hasilnya secara tertulis dan lisan.

\section{Tujuan}
Tujuan pembuatan skripsi adalah:
\begin{enumerate}
\item Agar mahasiswa memiliki kemampuan, keterampilan, dan pengalaman untuk memecahkan salah satu masalah dalam bidang rekayasa sistem komputer.
\item Memberi kesempatan kepada mahasiswa untuk dapat memahami dan mendalami salah satu cabang ilmu rekayasa sistem komputer.
\item Agar mahasiswa memiliki kemampuan untuk menyusun suatu karya ilmiah berdasarkan metode keilmuan.
\item Melatih mahasiswa menyajikan dan mempertanggungjawabkan hasil penelitiannya di hadapan penguji.
\end{enumerate}

\section{Ruang Lingkup dan Jenis Penelitian}
Topik skripsi harus berada dalam bidang keilmuan Rekayasa Sistem Komputer. Jenis penelitian yang dapat dijadikan skripsi antara lain:
\begin{enumerate}
\item perancangan dan implementasi sistem, baik perangkat keras, perangkat lunak, sistem tertanam (\textit{embedded system}), maupun sistem berbasis \textit{Internet of Things} (IoT) beserta pengujian kinerjanya;
\item penelitian eksperimental yang membandingkan metode, algoritma, atau konfigurasi sistem berdasarkan indikator kinerja yang terukur;
\item pengembangan atau penerapan algoritma dan model untuk menyelesaikan masalah tertentu, misalnya pengolahan sinyal, pengolahan citra, kendali, atau pembelajaran mesin;
\item studi kasus atau survei yang menganalisis kebutuhan, kinerja, atau keamanan suatu sistem pada objek nyata.
\end{enumerate}
Penelitian dapat bersifat kuantitatif, kualitatif, atau gabungan keduanya sesuai rumusan masalah dan tujuan penelitian.

\section{Susunan Skripsi}
Skripsi terdiri atas tiga bagian, yaitu bagian awal, bagian isi, dan bagian akhir. Susunan lembarannya, mulai dari halaman sampul sampai lampiran, berturut-turut diatur sebagai berikut:
\begin{enumerate}
\item Halaman sampul depan
\item Lembar identitas dan persetujuan pembimbing
\item Halaman pengesahan
\item Pernyataan keaslian
\item Abstrak (bahasa Indonesia) dan \textit{Abstract} (bahasa Inggris)
\item Kata pengantar
\item Daftar isi
\item Daftar tabel
\item Daftar gambar
\item Daftar lampiran
\item BAB I Pendahuluan
\item BAB II Tinjauan Pustaka
\item BAB III Metode Penelitian
\item BAB IV Hasil Penelitian dan Pembahasan
\item BAB V Kesimpulan dan Saran
\item Daftar pustaka
\item Lampiran
\end{enumerate}
Uraian setiap bagian dijelaskan pada Bab II, aturan penulisannya pada Bab III, dan tata cara pelaksanaannya pada Bab IV.

\section{Penyusunan Usulan Penelitian}
Sebelum membuat suatu usulan penelitian, perlu diingat bahwa antara judul penelitian, latar belakang, masalah, tujuan dan kegunaan, kerangka pikir, hipotesis, serta materi dan metode penelitian yang direncanakan harus runtut yang bertujuan untuk memudahkan evaluasi tentang dapat atau tidaknya penelitian dilakukan. Usulan penelitian harus sudah disetujui oleh dosen pembimbing sebelum penelitian dilaksanakan.

\subsection{Halaman Sampul}
Contoh halaman sampul usulan penelitian sesuai dengan halaman sampul Pedoman Penulisan Skripsi Universitas Muhammadiyah Buton.

\subsection{Halaman Pengajuan}
Halaman Pengajuan memuat judul penelitian, kata-kata Usulan Penelitian, Oleh, nama mahasiswa, dan NPM, sesuai dengan Halaman Pengajuan Pedoman Penulisan Skripsi Universitas Muhammadiyah Buton.

\subsection{Halaman Pengesahan}
Halaman ini memuat judul penelitian, kata-kata Usulan Penelitian, Oleh, nama mahasiswa, NPM, kolom tanda tangan persetujuan para dosen pembimbing, Ketua Program Studi dan Dekan, sesuai dengan Halaman Pengesahan Pedoman Penulisan Skripsi Universitas Muhammadiyah Buton.

\subsection{Pendahuluan}
Pendahuluan disusun secara sistematis dan terarah dalam urutan logika yang memberi justifikasi terhadap kerangka pikir, latar belakang masalah, tujuan, kegunaan, rumusan masalah dan hipotesis penelitian yang diajukan.

Pendahuluan harus menguraikan secara jelas dan ringkas tentang latar belakang masalah yang dihadapi dan perlunya penelitian dilakukan untuk memecahkan masalah yang dihadapi. Uraian tersebut dapat dibuat berdasarkan kepustakaan, kenyataan di lapang atau kombinasi keduanya.

Tujuan dan manfaat penelitian harus pula dipaparkan secara jelas tetapi singkat. Bab Pendahuluan ini juga dapat memuat hipotesis.

\subsection{Latar Belakang}
Latar belakang berisi justifikasi terhadap penelitian yang akan dilakukan dalam kaitannya dengan tujuan penelitian yang diajukan, uraian yang mengemukakan permasalahan-permasalahan yang secara umum dihadapi dan terkait dengan judul penelitian. Pendapat peneliti baik yang ditulis dalam bentuk \textit{textbook} maupun jurnal hasil penelitian yang relevan dengan judul penelitian yang diusulkan perlu juga dikemukakan.

\subsection{Rumusan Masalah}
Rumusan masalah didasarkan pada latar belakang yang telah dikemukakan. Dari latar belakang masalah disusun satu atau lebih rumusan masalah. Bentuk rumusan masalah umumnya berupa pertanyaan, dapat pula berbentuk pernyataan yang memerlukan jawaban melalui penelitian ilmiah.

\subsection{Tujuan Penelitian}
Tujuan penelitian disusun berdasarkan rumusan masalah, sehingga perlu adanya konsistensi antara rumusan masalah dengan tujuan penelitian. Tujuan penelitian menyatakan segala sesuatu yang ingin diperoleh dan diketahui melalui penelitian. Tujuan penelitian berupa kalimat pernyataan singkat, jelas dan spesifik tentang apa yang ingin diketahui dari kegiatan penelitian.

\subsection{Manfaat Penelitian}
Manfaat penelitian menguraikan tentang manfaat yang diperoleh apabila tujuan penelitian tercapai dan difokuskan pada pengembangan keilmuan, informasi bagi pembuat kebijakan dan umpan balik bagi para pengguna terutama dalam pembangunan di bidang rekayasa sistem komputer.

\subsection{Kerangka Pikir}
Kerangka pikir menguraikan tentang hasil-hasil penelitian atau pendapat terdahulu yang serupa, yang dapat digunakan untuk landasan berpijak dalam menentukan pengajuan hipotesis. Kerangka pemikiran merupakan justifikasi ilmiah terhadap penelitian yang akan dilakukan. Pada hakekatnya kerangka pikir harus didukung oleh landasan teoritis dan ditunjang oleh landasan informasi yang bersumber dari berbagai laporan ilmiah, sehingga menghasilkan suatu pemikiran baru dalam pendekatannya. Pemikiran baru tersebut dapat merupakan modifikasi atau penyempurnaan, gabungan atau perombakan total terhadap pendekatan lama yang tercermin dalam hipotesis penelitian.

\section{Prinsip Penulisan Ilmiah}
Dalam menyusun skripsi, mahasiswa wajib memegang prinsip berikut:
\begin{enumerate}
\item \textbf{Keaslian.} Skripsi adalah karya sendiri. Penggunaan gagasan, data, gambar, atau kalimat orang lain harus disertai sitasi dan dicantumkan dalam daftar pustaka. Plagiarisme dalam bentuk apa pun tidak dibenarkan dan berakibat pembatalan skripsi.
\item \textbf{Keruntutan.} Judul, latar belakang, rumusan masalah, tujuan, metode, hasil, dan kesimpulan harus saling berkaitan dan konsisten.
\item \textbf{Keterujian.} Metode dan data diuraikan cukup rinci sehingga orang lain dapat mengulangi penelitian dengan cara yang sama.
\item \textbf{Kejujuran data.} Data tidak boleh direkayasa, dihilangkan secara selektif, atau diubah untuk menyesuaikan hasil yang diharapkan.
\item \textbf{Bahasa yang baik.} Skripsi ditulis dengan bahasa Indonesia yang baku, singkat, jelas, dan efektif.
\end{enumerate}
